\section{特权机器完整走读：从用户指令到内核返回}

现在把前面的硬件拼起来，走读一台带 OS 的机器如何处理用户程序。目标不是背 x86-64 细节，而是建立完整路径：用户指令、异常、中断、系统调用、页表、设备、调度如何互相连接。

\subsection{机器启动后的两类代码}

系统启动后，内核代码运行在特权模式，用户程序运行在非特权模式。它们不是简单“谁更重要”，而是 CPU 对它们执行不同规则。

\begin{longtable}{p{0.22\linewidth}p{0.34\linewidth}p{0.32\linewidth}}
\toprule
能力 & 用户态 & 内核态 \\
\midrule
普通算术和分支 & 可以 & 可以 \\
访问用户页 & 可以，受页表限制 & 可以，但必须小心用户指针 \\
访问内核页 & 不可以 & 可以 \\
修改页表根 & 不可以 & 可以 \\
开关中断 & 不可以 & 可以 \\
访问设备 MMIO & 通常不可以 & 驱动可以 \\
返回任意模式 & 不可以 & 通过受控返回路径 \\
\bottomrule
\end{longtable}

用户态的限制不是库函数实现的，而是 CPU 和页表强制的。用户程序绕过 libc 直接写设备地址，也会 fault。

\subsection{路径一：普通用户指令}

用户程序执行普通加法：

\begin{lstlisting}
add rax, rbx
\end{lstlisting}

硬件动作是取指、译码、读寄存器、ALU 加法、写回。没有进入内核。OS 不参与每一条普通指令，否则程序会慢到不可用。OS 只通过页表、特权级、调度器和中断在边界处介入。

\subsection{路径二：用户 load 命中}

用户程序读取已映射内存：

\begin{lstlisting}
mov rax, qword ptr [rbx]
\end{lstlisting}

硬件查询 TLB，确认页表权限允许用户读，cache 命中后返回数据。整个过程不进内核。虚拟内存能高效运行，是因为权限检查由硬件快速完成，而不是内核逐条检查。

\subsection{路径三：用户 load 缺页}

若 TLB/page table 发现 PTE 不 present，CPU 进入 page fault。

\begin{lstlisting}
hardware:
  save faulting RIP
  save fault address and error code
  switch to kernel entry

kernel:
  find VMA
  if legal:
    allocate/load page
    install PTE
    return to same instruction
  else:
    deliver SIGSEGV or kill process
\end{lstlisting}

这条路径说明 page fault 是硬件和 OS 的合作。硬件只知道访问无法完成；OS 判断这是合法懒分配、文件映射、COW，还是非法地址。

\subsection{路径四：用户写只读 COW 页}

\code{fork} 后，父子进程可共享同一物理页，但 PTE 标为只读。任何一方写入时触发保护异常。

\begin{lstlisting}
write to COW page:
  hardware sees PTE writable = 0
  page fault with write error
  kernel sees VMA allows write and page is COW
  allocate new page
  copy old content
  update PTE writable to new page
  return and retry store
\end{lstlisting}

没有页表权限位和精确异常，COW 就无法透明实现。用户程序以为自己只是普通写内存，OS 在异常路径完成复制。

\subsection{路径五：系统调用}

用户程序要写文件：

\begin{lstlisting}
write(fd, buf, len)
\end{lstlisting}

libc wrapper 把系统调用号和参数放进寄存器，执行 syscall 指令。CPU 切换到内核入口。内核保存寄存器，检查 fd、用户 buffer、权限和对象状态，然后执行文件路径。

\begin{lstlisting}
sys_write:
  file = current.fd_table[fd]
  check file is writable
  copy bytes from user buffer
  write into page cache or device path
  return bytes_written or -errno
\end{lstlisting}

这里有两个边界。系统调用返回成功表示内核接受并完成了该调用语义，不一定表示数据持久化。用户 buffer 复制成功表示当时那段用户地址可读，不表示用户以后不会修改原 buffer。

\subsection{路径六：定时器中断打断用户程序}

用户程序正在运行，定时器到了。

\begin{lstlisting}
hardware:
  finish current instruction boundary
  save user RIP/RFLAGS
  enter timer vector

kernel timer handler:
  account time to current task
  run expired timers
  if time slice exhausted:
    set need_resched

return path:
  if need_resched:
    schedule()
  return to chosen task
\end{lstlisting}

若调度器选择另一个任务，返回用户态时就不回到原任务，而是恢复新任务的 RIP/RSP/寄存器/页表。用户程序感受到的是“自己某段时间没运行”，不是显式函数调用。

\subsection{路径七：设备完成中断}

磁盘或网卡完成 DMA 后发中断。

\begin{lstlisting}
device:
  writes data/status to memory
  raises MSI-X interrupt

kernel:
  enter IRQ handler
  acknowledge interrupt
  read completion queue
  unmap DMA buffer
  mark request done
  wake waiting task
\end{lstlisting}

这条路径说明设备和 CPU 是并发的。CPU 提交请求后可能去运行别的任务；设备完成时异步打断某个 CPU；内核把 completion 转换成等待队列唤醒或上层回调。

\subsection{路径八：返回用户态前的检查}

内核不能处理完就随便返回。返回用户态前，要检查信号、抢占、调试、审计、单步、恢复寄存器和返回模式。

\begin{lstlisting}
prepare_return_to_user:
  if signal_pending:
    build signal frame on user stack
  if need_resched:
    schedule
  run syscall exit tracing/audit
  restore user registers
  return with user mode bit
\end{lstlisting}

返回路径是安全边界。若返回模式错、栈错、寄存器泄露、信号帧构造错，都可能造成崩溃或漏洞。很多内核 entry 代码复杂，是因为它同时处理性能快路径和安全慢路径。

\subsection{一张总图：OS 介入点}

\begin{longtable}{p{0.22\linewidth}p{0.34\linewidth}p{0.32\linewidth}}
\toprule
事件 & 谁先发现 & 内核做什么 \\
\midrule
普通加法 & 不进入内核 & 不介入 \\
合法内存命中 & MMU/TLB/cache & 不介入 \\
缺页 & MMU & 分配/装载页面或报错 \\
权限错误 & MMU & COW、信号或杀进程 \\
系统调用 & 用户执行陷入指令 & 检查对象和权限，执行请求 \\
定时器 & timer hardware & 统计时间、调度、超时 \\
设备完成 & device/interrupt controller & completion、唤醒、错误处理 \\
非法指令 & decoder/privilege check & 信号、模拟或杀进程 \\
\bottomrule
\end{longtable}

这张表是后续全书的入口图。OS 不控制每个周期，但它控制关键边界；硬件不懂进程和文件，但它能在边界处把控制权交给内核。

\subsection{本节验收}

\begin{rubricbox}
\begin{itemize}
\tightlist
\item 能说明普通指令为什么不进入内核，系统仍然能被 OS 管理。
\item 能画出合法 load、缺页 load、COW store 三条路径的差异。
\item 能从 \code{write(fd, buf, len)} 说到寄存器 ABI、系统调用入口、fd 表、用户复制、page cache。
\item 能解释定时器中断如何让当前任务被换出。
\item 能解释设备 completion 如何唤醒等待 I/O 的任务。
\item 能说明返回用户态前为什么要处理信号、抢占和 tracing。
\end{itemize}
\end{rubricbox}
